% =====================================================================
%  Interactive Storybook - Project Documentation
%  Module: Natural User Interfaces (NUI), HTW Berlin
%
%  HOW TO USE THIS FILE ON OVERLEAF
%  1. Create a new project -> "Upload Project" and drop this .tex in,
%     OR "New Project" -> "Blank" and paste this file's contents.
%  2. Menu (top-left) -> set the compiler to "pdfLaTeX".
%  3. The bibliography is embedded at the very top via filecontents,
%     so no separate .bib upload is needed. It uses biblatex + biber,
%     which Overleaf runs automatically. Just press "Recompile".
%  4. Upload your images into an "images/" folder and replace each
%     \imgplaceholder{...} with the \includegraphics line shown in the
%     comment beside it.
%  5. Search for "TODO" to find everything you still need to fill in
%     (author names, matriculation numbers, per-member reflections,
%     photos, exact numbers you want to double-check).
% =====================================================================

% ---- Embedded bibliography (writes references.bib on first compile) ----
\begin{filecontents}[overwrite]{references.bib}
@inproceedings{ishii1997tangible,
  author    = {Ishii, Hiroshi and Ullmer, Brygg},
  title     = {Tangible Bits: Towards Seamless Interfaces between People,
               Bits and Atoms},
  booktitle = {Proceedings of the ACM SIGCHI Conference on Human Factors
               in Computing Systems (CHI '97)},
  pages     = {234--241},
  year      = {1997},
  publisher = {ACM}
}

@article{hornecker2006getting,
  author  = {Hornecker, Eva and Buur, Jacob},
  title   = {Getting a Grip on Tangible Interaction: A Framework on Physical
             Space and Social Interaction},
  journal = {Proceedings of the SIGCHI Conference on Human Factors in
             Computing Systems (CHI '06)},
  pages   = {437--446},
  year    = {2006},
  publisher = {ACM}
}

@inproceedings{nielsen1990heuristic,
  author    = {Nielsen, Jakob and Molich, Rolf},
  title     = {Heuristic Evaluation of User Interfaces},
  booktitle = {Proceedings of the SIGCHI Conference on Human Factors in
               Computing Systems (CHI '90)},
  pages     = {249--256},
  year      = {1990},
  publisher = {ACM}
}

@online{nielsen1994heuristics,
  author  = {Nielsen, Jakob},
  title   = {10 Usability Heuristics for User Interface Design},
  year    = {1994},
  note    = {Nielsen Norman Group},
  url     = {https://www.nngroup.com/articles/ten-usability-heuristics/}
}

@book{norman2013everyday,
  author    = {Norman, Donald A.},
  title     = {The Design of Everyday Things: Revised and Expanded Edition},
  year      = {2013},
  publisher = {Basic Books},
  address   = {New York}
}

@book{dourish2001where,
  author    = {Dourish, Paul},
  title     = {Where the Action Is: The Foundations of Embodied Interaction},
  year      = {2001},
  publisher = {MIT Press},
  address   = {Cambridge, MA}
}

@article{wilson2002six,
  author  = {Wilson, Margaret},
  title   = {Six Views of Embodied Cognition},
  journal = {Psychonomic Bulletin \& Review},
  volume  = {9},
  number  = {4},
  pages   = {625--636},
  year    = {2002}
}

@article{marshall2003tangible,
  author  = {Marshall, Paul and Price, Sara and Rogers, Yvonne},
  title   = {Conceptualising Tangibles to Support Learning},
  journal = {Proceedings of the Conference on Interaction Design and Children (IDC '03)},
  pages   = {101--109},
  year    = {2003},
  publisher = {ACM}
}

@online{nngroup_workbook,
  author = {{Nielsen Norman Group}},
  title  = {Heuristic Evaluation Workbook},
  year   = {2021},
  url    = {https://www.nngroup.com/articles/how-to-conduct-a-heuristic-evaluation/}
}
\end{filecontents}

\documentclass[11pt,a4paper]{article}

% ---- Encoding and fonts ----
\usepackage[T1]{fontenc}
\usepackage[utf8]{inputenc}
\usepackage{lmodern}

% ---- Layout ----
\usepackage[a4paper,margin=2.5cm]{geometry}
\usepackage{setspace}
\onehalfspacing
\usepackage{parskip}

% ---- Graphics, tables, colour ----
\usepackage{graphicx}
\usepackage{booktabs}
\usepackage{array}
\usepackage{caption}
\usepackage{xcolor}
\definecolor{accent}{HTML}{B8306A}
\definecolor{codebg}{HTML}{F4F6F8}

% ---- Code listings ----
\usepackage{listings}
\lstset{
  basicstyle=\ttfamily\small,
  backgroundcolor=\color{codebg},
  frame=single,
  rulecolor=\color{gray!40},
  breaklines=true,
  columns=fullflexible,
  showstringspaces=false
}

% ---- Links and bibliography ----
\usepackage[backend=biber,style=numeric,sorting=none]{biblatex}
\addbibresource{references.bib}
\usepackage[hidelinks]{hyperref}

% ---- Helper: image placeholder box (compiles before you add images) ----
% Replace the whole \imgplaceholder{...} with, e.g.:
%   \includegraphics[width=0.8\textwidth]{images/system-diagram.png}
\newcommand{\imgplaceholder}[1]{%
  \begin{center}
    \setlength{\fboxsep}{12pt}%
    \fbox{\parbox[c][3.2cm][c]{0.82\textwidth}{\centering\itshape #1}}
  \end{center}%
}

% =====================================================================
\begin{document}

% ------------------------- TITLE PAGE --------------------------------
\begin{titlepage}
\centering
\vspace*{2cm}

{\Large Hochschule f\"ur Technik und Wirtschaft Berlin\par}
{\large Module: Natural User Interfaces\par}
\vspace{2.5cm}

{\Huge\bfseries Interactive Storybook\par}
\vspace{0.5cm}
{\Large A Tangible NFC Interface for Children\par}
\vspace{2.5cm}

{\large Project Documentation\par}
\vspace{2cm}

% TODO: fill in your group members and matriculation numbers.
\begin{tabular}{ll}
Author 1 & Matriculation No.\\
Author 2 & Matriculation No.\\
Author 3 & Matriculation No.\\
Author 4 & Matriculation No.\\
\end{tabular}

\vfill
% TODO: set the correct term / supervisor.
{Supervisor: Paulina Stefanovic \quad\textbar\quad Winter Term 2025/26\par}
\end{titlepage}

% ------------------------- TABLE OF CONTENTS -------------------------
\tableofcontents
\newpage

% =====================================================================
\section{Introduction}
% =====================================================================

\subsection{Topic}
This project presents the \emph{Interactive Storybook}, an offline,
tangible storytelling game for children. Instead of clicking buttons on a
screen, a child advances a branching narrative by placing physical
\emph{NFC cards} on a reader. Each card carries a symbolic meaning --- a
story to begin, an action to take, or a system command --- and the software
responds by moving the story forward, updating a small inventory, and
displaying the next illustrated scene.

The system combines an Arduino RC522 NFC reader with a Python application.
The reader transmits the unique identifier (UID) of each scanned card over
USB, and the application maps that identifier to a symbolic card, feeds it
into a story state machine, and renders the resulting scene in a full-screen,
child-friendly interface.

\subsection{Goals}
The project pursues three primary goals:
\begin{itemize}
  \item \textbf{Tangible interaction.} The physical card is the only means of
        input. The screen reports state but offers no gameplay controls, so
        the interaction is genuinely embodied rather than a graphical
        interface with unusual peripherals.
  \item \textbf{Accessibility for pre-readers.} The target audience cannot
        yet navigate menus or read long instructions. Interaction is driven
        by recognisable icons and colours instead of text.
  \item \textbf{Content without code.} Stories are stored as external data
        files. Adding a new story or card requires no change to the program
        logic, which keeps the system maintainable and extensible.
\end{itemize}

\subsection{Motivation and Target Audience}
Conventional screen interfaces demand skills a young child has not yet
developed: reading, pointer control, and menu navigation. A tangible
interface removes these barriers by letting the child interact with a
physical object they already understand how to handle --- picking it up and
putting it down \cite{ishii1997tangible}. This aligns the project directly
with the concerns of a Natural User Interface: an interface that feels
natural because it builds on existing human capabilities rather than
learned conventions.

The intended user is a child aged roughly six to nine, a pre-reader or early
reader, typically accompanied by a parent or educator. The context of use is
a shared table: the screen stands upright, the reader sits in front of it,
and a deck of physical cards is spread within reach. The setting is
deliberately social --- the cards are visible to everyone at the table, which
turns each decision into a shared, negotiable object rather than a private
action on a personal screen.

% =====================================================================
\section{Theoretical Background}
% =====================================================================

\subsection{Natural User Interfaces and Tangible Interaction}
A Natural User Interface (NUI) aims to make interaction feel natural by
leveraging skills users already possess. Tangible User Interfaces (TUIs) are
a prominent branch of this idea. Ishii and Ullmer's seminal work on
\emph{Tangible Bits} argues for coupling digital information to physical
objects and surfaces, so that the interface becomes graspable and inhabits
the physical world rather than being confined behind glass
\cite{ishii1997tangible}. Our project is a direct application of this
principle: the digital story is coupled to a set of physical cards, and the
act of placing a card is the act of making a choice.

Hornecker and Buur provide a framework for analysing tangible interaction
that goes beyond the object itself to include physical space and social
interaction \cite{hornecker2006getting}. Their emphasis on the social and
spatial dimension is directly relevant here: because the cards lie openly on
a shared table, the design supports collaboration and turn-taking between
children, a property a single-pointer graphical interface cannot easily
provide.

\subsection{Embodied Cognition}
The theoretical grounding for ``natural'' interaction is embodied cognition,
the view that cognition is shaped by the body's interaction with the world
rather than occurring purely in the abstract \cite{wilson2002six}. Dourish
develops this into a foundation for interaction design, arguing that meaning
arises through embodied, situated action \cite{dourish2001where}. For a young
child, whose thinking is strongly tied to physical manipulation, an interface
that is operated by hand is not merely convenient but cognitively
appropriate. This motivated our decision to explore the concept physically
through bodystorming before committing to a screen design.

\subsection{Affordances and Signifiers}
Norman's account of affordances and signifiers frames how users perceive what
actions are possible \cite{norman2013everyday}. A physical card affords
picking up and placing; a printed icon signifies its meaning. Norman's three
guiding questions --- \emph{How do I know an action is possible? How do I do
it? How do I know I succeeded?} --- structured our interface decisions, and
recur in our evaluation: for example, the persistent ``last scanned card''
readout exists specifically to answer the third question.

\subsection{Usability Evaluation Methods}
To evaluate the design we used heuristic evaluation, an inspection method in
which evaluators judge an interface against a set of recognised usability
principles \cite{nielsen1990heuristic}. We applied Nielsen's ten usability
heuristics \cite{nielsen1994heuristics} using the Nielsen Norman Group
workbook \cite{nngroup_workbook}. Heuristic evaluation was appropriate for
our stage of development because it identifies a large share of usability
problems with a small number of evaluators and can be conducted on a
low-fidelity prototype before implementation is complete.

\subsection{Related Work}
Tangible interfaces for children's learning are an established research area.
Marshall, Price, and Rogers discuss how tangibles can support learning
through physical engagement \cite{marshall2003tangible}. Our project sits
within this tradition but focuses on branching narrative and moral choice
rather than abstract concept learning, using the physical card as the unit
of narrative decision.

% =====================================================================
\section{Concept}
% =====================================================================

\subsection{Requirements}
From the goals and target audience we derived the following requirements:
\begin{itemize}
  \item Input exclusively through physical NFC cards; no gameplay buttons on
        screen.
  \item Full offline operation --- no network, no accounts, no cloud services.
  \item Icon- and colour-driven interface legible to a pre-reader.
  \item Branching stories with multiple endings and a simple inventory.
  \item Robustness suitable for unsupervised use: invalid input must never
        break the story or crash the application.
  \item New stories and cards addable through data files without code
        changes.
\end{itemize}

\subsection{User Stories}
% TODO: this is a condensed version. The full set of six user stories is in
% docs/USER_STORIES.md -- expand here or keep concise, your choice.
The design was framed around concrete user stories. A representative example:

\begin{quote}
\emph{As a} six-to-nine-year-old child at a shared table, \emph{I want to}
choose what the character does by placing a picture card on the reader,
\emph{so that} I can steer the story myself without reading menus or using a
mouse.
\end{quote}

Further stories covered starting a story, making a choice (the core loop),
recovering from an invalid card, encountering a path that requires a
collected item, replaying to explore other branches, and two children
sharing one table. Each story pairs a physical \emph{trigger action}
(placing a card) with a defined system response and an articulated design
value.

\subsection{Prototyping and UX Methods}

\subsubsection{Bodystorming}
Because our system is a tangible interface, we explored the interaction
physically before designing any screen, using bodystorming --- physically
acting out the scenario with paper props to surface assumptions about how a
child holds, places, and reacts to a card.

% TODO: replace with a real photo from your bodystorming session.
\imgplaceholder{Figure placeholder: photo(s) of the bodystorming session ---
the paper cards, the cardboard ``reader'', and a participant acting out a
scan. Replace with \texttt{\textbackslash includegraphics}.}
\captionof{figure}{Bodystorming with paper prototypes.}

% TODO: write 3-5 sentences on what the bodystorming actually revealed.
\textbf{Findings (to complete from your session).} Note here the concrete
insights: for example, whether participants tapped instead of placing the
card, how long they held it, and whether they knew a scan had registered.
These observations feed directly into the interface decisions below.

\subsubsection{Low-Fidelity Prototype and Wireframe}
We then produced a low-fidelity wireframe. Because the interface is not only
the screen but the whole table, the wireframe has two layers: the physical
setup (screen, reader, and card deck seen from above) and the three screen
states (start, scene, ending) with a persistent footer.

% TODO: replace with your paper prototype / wireframe photos or exports.
\imgplaceholder{Figure placeholder: the two-layer wireframe --- (a) top-down
physical setup, (b) the three screen states with the persistent footer.}
\captionof{figure}{Two-layer wireframe: physical setup and screen states.}

\subsubsection{Heuristic Evaluation}
We evaluated the wireframe and early prototype against Nielsen's ten
heuristics \cite{nielsen1994heuristics}. Selected findings are summarised in
Table~\ref{tab:heuristics}; the full workbook is included in the appendix
material.

\begin{table}[htbp]
\centering
\caption{Selected heuristic evaluation findings and fixes.}
\label{tab:heuristics}
\small
\begin{tabular}{@{}p{3.1cm}p{5.3cm}p{5.3cm}@{}}
\toprule
\textbf{Heuristic} & \textbf{Issue found} & \textbf{Fix applied} \\
\midrule
Consistency \& standards &
Choice colours were assigned by on-screen position, so the same card appeared
in a different colour in each scene. &
Every card now has one fixed colour and icon, matched between the printed
card and the screen. \\
\addlinespace
User control \& freedom &
A child who picked the wrong story had no way back; the window was fullscreen
with no exit. &
Added a physical ``Home'' card and an on-screen Home control as an emergency
exit; bound Escape to leave fullscreen. \\
\addlinespace
Visibility of system status &
An invalid scan was shown only as small footer text for four seconds, easily
missed. &
Identified as a priority fix: central feedback and sound (see
limitations). \\
\addlinespace
Help \& documentation &
Nothing told a first-time child that cards are placed, not tapped. &
Start-screen hint and printed instruction card proposed. \\
\bottomrule
\end{tabular}
\end{table}

% =====================================================================
\section{Implementation}
% =====================================================================

\subsection{Technologies Used}
The system is deliberately built from a minimal, well-established technology
stack so that it installs with a single command and runs offline on any
machine in the room.

\begin{table}[htbp]
\centering
\caption{Technologies and their role in the system.}
\label{tab:tech}
\small
\begin{tabular}{@{}llp{7.2cm}@{}}
\toprule
\textbf{Technology} & \textbf{Type} & \textbf{Role} \\
\midrule
Arduino Uno + RC522 & Hardware & Reads NFC card UIDs and sends them over USB
serial. \\
MFRC522 library & Firmware & Drives the RC522 reader over SPI. \\
Python 3.12 & Language & Application logic. \\
Tkinter & Std.\ library & Full-screen graphical user interface. \\
pyserial & Third-party & Serial communication with the Arduino. \\
Pillow (PIL) & Third-party & Loading, scaling, and placeholder generation of
scene images. \\
pytest & Third-party & Automated test suite (development only). \\
\bottomrule
\end{tabular}
\end{table}

A notable property is that only three third-party libraries are required;
everything else, including the entire user interface, uses the Python
standard library.

\subsection{System Architecture}
The application follows a strict four-layer architecture. The two layers that
touch the outside world (serial I/O and the graphical interface) sit at the
edges, so the core game logic is independent of both hardware and GUI and can
be tested without either.

% TODO: export the architecture diagram to an image and include it here.
% A ready diagram exists in docs/ARCHITECTURE_AS_BUILT.md (Mermaid), and the
% story maps in docs/img/ were rendered the same way.
\imgplaceholder{Figure placeholder: SYSTEM DIAGRAM --- schematic of the
architecture showing hardware and software components and the data exchange
between them (NFC card $\rightarrow$ RC522 $\rightarrow$ Arduino $\rightarrow$
USB serial $\rightarrow$ SerialReader $\rightarrow$ CardManager $\rightarrow$
StoryEngine $\rightarrow$ GameUI). Replace with
\texttt{\textbackslash includegraphics}.}
\captionof{figure}{System architecture and data flow.}

The data flow proceeds as follows. A child places a card on the reader; the
Arduino reads its UID and sends it as a line of text over USB serial at
115200 baud. On the Python side, a background thread receives the UID,
validates and de-bounces it, and hands it to the main application thread. The
\texttt{CardManager} maps the UID to a symbolic card; the \texttt{StoryEngine}
validates the card against the current scene and updates the story state; and
the \texttt{GameUI} renders the resulting scene. Because the serial reader
runs on a separate thread while Tkinter is not thread-safe, all state changes
are marshalled onto the main thread; this threading boundary is the single
most important correctness detail in the system.

\subsection{The Story Engine}
The core of the application is the story engine, a state machine that tracks
the active story, the current scene, the player's inventory, and whether an
ending has been reached. Every scanned card enters through a single method
and leaves as a typed result describing what happened. Stories are stored as
external JSON files: each scene defines its text, image, the choices
available (mapping a card name to a target scene), and optional item effects.
Adding a story therefore requires only a new data file, not a code change.

A subtle but important design rule is that items are \emph{not} collected by
scanning a card. Instead, a scene grants an item when the child enters it,
and item requirements are checked on the \emph{target} scene before a move is
committed. This means the physical card cannot ``cheat'' a locked path: the
child must actually have visited the scene that grants the required item.

\subsection{Hardware Integration}
The Arduino firmware is intentionally minimal: it reads a card UID and
transmits it, and contains no game logic whatsoever. This separation keeps
the microcontroller simple and, more importantly, makes the reader
replaceable --- the game would work with any device that can send a UID over
serial. The Python side auto-detects the serial port and reconnects
automatically if the reader is unplugged, so the application survives a
disconnection during use without losing story state.

\subsection{The Debug Mode (Testing Without Hardware)}
During development we implemented a debug mode inside the application. Instead
of scanning physical NFC cards, on-screen buttons and keyboard shortcuts
simulate each card. A simulated scan is not a separate code path: it enters
the same pipeline one step later, looking the card up by name instead of by
UID, so everything after that point --- validation, item gates, rendering ---
is the identical code that runs with the hardware.

This had two benefits. It let us test story branches, transitions, and error
handling quickly without the Arduino; and it meant three group members could
work on stories, interface, and logic while one member had the physical
reader. The debug mode was therefore as much an organisational decision as a
technical one, and it is the main reason the architecture keeps the core
logic independent of the hardware.

\subsection{Significant Milestones and Turning Points}
\begin{enumerate}
  \item \textbf{The Story Engine.} The branching state machine was built and
        tested first, with stories as data, before any hardware existed.
  \item \textbf{Debug mode.} On-screen simulation of card scans made the
        hardware optional and enabled parallel work across the group.
  \item \textbf{Arduino RC522 integration.} Reliable serial communication
        made the project genuinely interactive, with deliberately minimal
        firmware.
  \item \textbf{Turning point: the card as sole input.} Gameplay controls
        were removed from the screen entirely; clicking survives only as a
        development tool. This was a deliberate constraint --- an on-screen
        button would be faster than finding a card, and would reduce the
        tangible interaction to decoration.
  \item \textbf{Iterative UI redesign.} Several rounds toward a brighter,
        child-friendly look; the most significant change --- one fixed colour
        and icon per card, matched to the printed card --- came from
        usability testing rather than aesthetics.
  \item \textbf{From prototype to system.} Printing the story, action, and
        system cards marked the transition from a software prototype to a
        complete interactive system testable with real users.
\end{enumerate}

\subsection{The Final System}
% TODO: add screenshots of the running application (start / scene / ending)
% and photos of the finished physical cards and a user interacting with it.
\imgplaceholder{Figure placeholder: screenshots of the running application ---
start screen, a story scene, and an ending screen.}
\captionof{figure}{The final interface: start, scene, and ending screens.}

\imgplaceholder{Figure placeholder: the printed card deck and a child (or
stand-in) placing a card on the reader --- the end prototype in use.}
\captionof{figure}{The end prototype in use.}

The finished system contains three stories with 75 scenes and 8 endings in
total, driven by a deck of thirteen cards (three story cards, seven action
cards, one item card, and two system cards). The application logic is
covered by an automated test suite of over one hundred tests that require no
hardware to run.

% =====================================================================
\section{Conclusion}
% =====================================================================

\subsection{Summary}
We designed and built a tangible NFC storytelling system for children,
grounded in the principles of tangible interaction and embodied cognition.
The physical card is the sole means of input, the interface is legible to a
pre-reader, and the branching content is fully data-driven. The design was
shaped by bodystorming and refined through heuristic evaluation, and several
concrete usability problems were identified and fixed as a result.

\subsection{Limitations}
Several limitations arise from technology and target audience:
\begin{itemize}
  \item \textbf{Feedback for invalid input.} An invalid scan is currently
        signalled only by brief on-screen text, which a child looking down at
        the cards may miss. Auditory feedback is the clear next step.
  \item \textbf{Reader range.} The RC522 has a short range and can re-read a
        card left near it; this is mitigated in software but is best handled
        physically by moving played cards away from the reader.
  \item \textbf{Evaluation scope.} Our evaluation used heuristic inspection
        and bodystorming. Testing with a larger group of children in the
        target age range remains the most important validity limitation and
        the clearest direction for future work.
  \item \textbf{No persistence.} A session is designed for a single sitting;
        there is no save-and-resume feature.
\end{itemize}

\subsection{Outlook}
Future work includes auditory feedback and narration (which would also make
the system usable by children who cannot read at all), a larger library of
stories, and a structured usability study with children in the target group.
The data-driven architecture means new stories can be added without touching
the code, which makes content expansion straightforward.

\subsection{Self-Reflection}
% TODO: EACH group member writes a short paragraph. Cover: your work share,
% what you learned, and hurdles you overcame. Keep it honest and specific.

\textbf{Member 1 (Name).}
% Work share: ... Knowledge gained: ... Hurdles: ...

\textbf{Member 2 (Name).}
%

\textbf{Member 3 (Name).}
%

\textbf{Member 4 (Name).}
%

% =====================================================================
\printbibliography[title={Bibliography}]
\addcontentsline{toc}{section}{Bibliography}
% =====================================================================

\end{document}
